\documentclass{resume} % Use the custom resume.cls style

\usepackage[left=0.4in,top=0.4in,right=0.4in,bottom=0.4in]{geometry}
\usepackage[utf8]{inputenc}
\usepackage[T2A]{fontenc}
\usepackage[russian]{babel}
\usepackage{fontawesome5}
\usepackage{graphicx}
\usepackage{tabularx}

\name{Караганов Павел}
\address{Python Backend Разработчик (Стажёр)}
\address{
    \href{mailto:karaganovpavel119@yandex.ru}{\faEnvelope\hspace{1}karaganovpavel119@yandex.ru} \hspace{2ex}
    \href{https://t.me/Antichnoste}{\faTelegram\hspace{1}@Antichnoste} \hspace{2ex}
    \href{https://github.com/Antichnoste}{\faGithub\hspace{1}github.com/Antichnoste}
}

\begin{document}

\begin{rSection}{О себе}
Студент 3 курса ИТМО (поступил по БВИ). Совмещаю сильную алгоритмическую базу (олимпиады по спортивному программированию) с практическим backend-опытом: проектирование REST API, работа с БД, DevOps и мониторинг. Самостоятельно на Python реализовал транслятор языка программирования и tick-точный симулятор процессора — проект, давший глубокое понимание внутреннего устройства языков и вычислительных систем. Целенаправленно развиваюсь в сторону Python backend, быстро обучаюсь и пишу эффективный, тестируемый код.
\end{rSection}

%----------------------------------------------------------------------------------------
%	EDUCATION SECTION
%----------------------------------------------------------------------------------------
\begin{rSection}{Образование}
{\bf Программная Инженерия}, бакалавриат ИТМО (Факультет ПИиКТ) \hfill {2024 - 2028}
\begin{itemize}
    \itemsep -3pt {} 
    \item Поступил без вступительных испытаний (БВИ) по результатам перечневых олимпиад.
    \item Профильные дисциплины: Алгоритмы и структуры данных, Математический анализ, Базы данных.
\end{itemize}
\end{rSection}

%----------------------------------------------------------------------------------------
% TECHINICAL STRENGTHS	
%----------------------------------------------------------------------------------------
\begin{rSection}{Навыки}
\begin{tabularx}{\linewidth}{ @{} >{\bfseries}l @{\hspace{4ex}} X }
Языки и Основы: & Python (OOP, стандартная библиотека, работа с бинарными данными), Java (Core, Collections, Multithreading, Stream API), SQL, C++ \\
Фреймворки и БД: & Spring (Boot, Security/JWT, Data), PostgreSQL, Liquibase \\
Архитектура и Сборка: & REST API, OpenAPI (Swagger), Maven \\
DevOps и Мониторинг: & Docker, Docker Compose, GitHub Actions (CI/CD), Prometheus, Grafana \\
Тестирование: & JUnit, Mockito, pytest, Yandex Tank
\end{tabularx}
\end{rSection}

%----------------------------------------------------------------------------------------
%	ACHIEVEMENTS SECTION
%----------------------------------------------------------------------------------------
\begin{rSection}{Достижения}
\begin{itemize}
    \itemsep -3pt {}
    \item \textbf{Спортивное программирование:} Призёр перечневых олимпиад, участник регионального этапа ВсОШ. Опыт оптимизации алгоритмов по времени и памяти.
    \item \textbf{Английский язык:} B2 (Upper-Intermediate).
\end{itemize}
\end{rSection} 

%----------------------------------------------------------------------------------------
%	WORK EXPERIENCE / PROJECTS SECTION
%----------------------------------------------------------------------------------------
\begin{rSection}{Практический опыт и Проекты}

\item \textbf{\href{https://github.com/Antichnoste/pipelined-lisp-cpu}{Разработка транслятора и симулятора процессора}} \hfill \textit{Python, Архитектура вычислительных систем}
\begin{itemize}
    \itemsep -3pt {}
    \item Спроектировал и реализовал на \textbf{Python} транслятор Lisp-подобного языка (S-expressions, рекурсивные функции) в бинарный машинный код для собственной процессорной архитектуры (Гарвардская, port-mapped I/O, null-terminated строки).
    \item Разработал tick-точный симулятор 3-стадийного конвейерного процессора с гибридным управлением (hardwired IF/ID + микрокод на стадии EX), поддерживающий stall/flush и условные переходы.
    \item Покрыл проект golden-тестами на \textbf{pytest} (8 сценариев: рекурсия, сортировка, ввод-вывод, арифметика) и настроил CI в \textbf{GitHub Actions}.
\end{itemize}

\item \textbf{\href{https://github.com/Antichnoste/observable-spatial-api}{Разработка клиент-серверного веб-приложения}} \hfill \textit{Java, Spring Boot, DevOps}
\begin{itemize}
    \itemsep -3pt {}
    \item Спроектировал REST API, настроил безопасность через \textbf{Spring Security (JWT)} и задокументировал контракты с помощью \textbf{OpenAPI (Swagger)}.
    \item Настроил интеграцию с базой данных \textbf{PostgreSQL} (управление миграциями через \textbf{Liquibase}) и упаковал инфраструктуру в \textbf{Docker Compose}.
    \item Для контроля надежности прикрутил сбор метрик приложения через \textbf{Prometheus}, настроил визуальные дашборды в \textbf{Grafana} и провел нагрузочное тестирование с помощью \textbf{Yandex Tank}.
    \item Внедрил пайплайн CI/CD в GitHub Actions (автоматическая сборка через \textbf{Maven}, тесты, линтеры).
\end{itemize}

\end{rSection} 

\end{document}